% ====================================================================
% INTRODUCCIÓN
% ====================================================================

\section{Introducción}

\subsection{Contexto y Motivación}

% Explicar el contexto del problema, por qué es importante,
% y qué motivó la selección de este proyecto
Entre 2021 y 2023 comenzo una fuerte contracción en la demanda de viviendas, especialmente en el segmento de casas nuevas sin subsidio que cayó un 11,8\% entre enero y octubre de 2023, las ventas de casas nuevas alcanzaron en octubre de este año a tan solo 330 unidades con una tasa de desistimientos del 33\% en casas sin subsidio y del 15\% en departamentos, niveles históricamente altos a esta alturas del año. Este desplome se debio al deterioro de las condiciones de crédito, el aumento de las tasas de interés, restricciones bancarias y el impacto de la inflación sobre la capacidad de endeudamiento manteniendose durante los años siguientes.
\\ 
Debido a lo anterior a comienzos del 2025 el alto stock de propiedades no vendidas, estimado en alrededor de 70.000 unidades, comenzo a transformar el mercado inmobiliario ofreciendo unas nueva posibilidades para compradores que buscan oportunidades de inversión o vivienda propia, especialmente mediante ofertas directas o grupos de compra \cite{biobio_ciper_2025}.
\\Segun el estudio ``monitor de viviendas 2025'' \cite{ipsos_monitor_2025} realizado por la consultora multinacional IPSOS en enero de este año 2025 mostro que el "60 \% de los encuestados estan contentos con la situacion actual de la vivienda" en Chile, de donde "tres de cada cuatro (74\%) propietarios esta satisfecho con su situacion de vivienda" y "solo cuatro de diez arrendatarios (39\%) dice estar feliz de su situacion habitacional", lo que alienta a nuevos clientes y a actuales arrendatarios a considerar la inversion inmobiliaria como un objetivo plausible en el actual contexto chileno.
\\
\\Respecto a esto se ha determinado que los principales desafios para los clientes al momento de comprar o arrendar una vivienda se identifican de la siguiente forma:

\begin{itemize}
    \item Precios altos de la propiedades
    \item Tasas de interes muy altas
    \item Aumento de los costos de construccion de viviendas y edificios
\end{itemize}
Esta dificultad para conseguir vivienda a llevado a que los clientes tomen en consideracion los siguientes aspectos al momentos de comprar una propiedad:

\begin{itemize}
    \item Relación precio-m²
    \item Tasa de delincuencia
    \item Ubicacion y servicios
    \item Acceso al transporte publico
    \item Espacio al aire libre
    \item Privacidad
    \item Buena infraestructura local
\end{itemize}

Por todo lo anterior, hoy se presenta una oportunidad interesante dentro del mundo inmobiliario chileno: la capacidad de influir en las decisiones de potenciales clientes para informar y convencer al cliente indeciso de comprar una vivienda con datos geoespaciales que le ayuden a determinar una opción de vivienda que satisfaga sus necesidades a un precio asequible.

\subsection{Objetivos}

\subsubsection{Objetivo General}

% Un objetivo general claro y alcanzable
Desarrollar un sistema de recomendación inmobiliaria inteligente que prediga la satisfacción residencial de propiedades en venta mediante la integración de características internas de las viviendas con factores geoespaciales del entorno urbano, empleando técnicas de Machine Learning (LightGBM) para proporcionar recomendaciones personalizadas que faciliten la toma de decisiones de compra en el mercado inmobiliario de la Región Metropolitana de Santiago.

\subsubsection{Objetivos Específicos}

\begin{enumerate}
    \item \textbf{Consolidar y normalizar información geoespacial}: Integrar 29 datasets geográficos heterogéneos de la Región Metropolitana, normalizando sistemas de coordenadas al CRS EPSG:32719 (UTM Zone 19S), validando calidad geométrica, y estandarizando estructuras de datos para garantizar consistencia y trazabilidad en el análisis espacial.
    
    \item \textbf{Cuantificar características espaciales del entorno urbano}: Generar una grilla sistemática de evaluación sobre el territorio de estudio y calcular 72 métricas espaciales cuantitativas por ubicación, incluyendo: (a) distancias euclidianas a 21 categorías de servicios urbanos, (b) densidades de servicios en radios de 300m, 600m y 1000m, y (c) índices compuestos de accesibilidad en 9 dimensiones de habitabilidad.
    
    \item \textbf{Recopilar y geocodificar datos del mercado inmobiliario}: Obtener mediante web scraping información de 7,702 propiedades en venta (casas y departamentos) de 4 comunas objetivo (La Reina, Ñuñoa, Santiago Centro, Estación Central), geocodificar direcciones mediante Google Maps API, y validar coherencia entre características de propiedades y ubicaciones geográficas.
    
    \item \textbf{Desarrollar proxy de satisfacción residencial}: Diseñar e implementar un índice de satisfacción (escala 1-10) que integre: (a) valor económico relativo (precio/m² contextualizado), (b) características físicas de la propiedad (espacio, distribución), y (c) accesibilidad a servicios urbanos, con capacidad de personalización mediante 5 perfiles de usuario (familia con niños, profesional joven, inversionista, adulto mayor, balanceado).
    
    \item \textbf{Entrenar y validar modelo predictivo de satisfacción}: Implementar modelo LightGBM para predecir satisfacción residencial utilizando 42 features (características internas + métricas espaciales), alcanzar $R^2 \geq 0.85$ en conjunto de prueba, validar mediante cross-validation 5-fold, analizar importancia de variables, y evaluar autocorrelación espacial de residuos para verificar ausencia de sesgo geográfico.

    \item \textbf{Desarrollar API de recomendación}: Crear sistema de predicción funcional que permita: (a) estimar satisfacción de propiedades nuevas, (b) generar rankings personalizados por perfil de usuario, (c) comparar múltiples opciones simultáneamente, y (d) explicar factores determinantes de cada predicción mediante análisis de importancia de variables.
    \item \textbf{Visualizar patrones espaciales y resultados}: Generar 3 mapas temáticos con elementos cartográficos estándar (ubicación de propiedades, distribución de precios, satisfacción predicha), 5 gráficos estadísticos de análisis exploratorio (histogramas, correlaciones, dispersión, importancia de variables), y 1 visualización interactiva web (mapa Folium con popups informativos y heatmap).
    \item \textbf{Evaluar desempeño comparativo de modelos}: Implementar línea base con Random Forest, comparar rendimiento con al menos 3 algoritmos adicionales (XGBoost, CatBoost, Neural Networks), documentar métricas de evaluación (R², RMSE, MAE, tiempo de entrenamiento), y justificar selección del modelo final mediante análisis cuantitativo de trade-offs precisión-interpretabilidad-eficiencia.
    
\end{enumerate}

\subsection{Preguntas de Investigación}

\
El proyecto busca responder las siguientes preguntas clave mediante análisis geoespacial cuantitativo:

\subsubsection{Pregunta Principal}

\textbf{¿Qué hace que una vivienda sea satisfactoria para vivir en el contexto urbano de Santiago?}

Esta pregunta se descompone en sub-preguntas específicas que el proyecto aborda:

\begin{enumerate}
    \item \textbf{Dimensión económica}: ¿Cuál es la relación entre el valor económico (precio/m²) y la satisfacción percibida? ¿Propiedades más caras son necesariamente más satisfactorias?
    
    \item \textbf{Dimensión espacial}: ¿En qué medida los factores externos (proximidad a servicios, áreas verdes, transporte) determinan la satisfacción residencial más allá de las características internas de la propiedad?
    
    \item \textbf{Heterogeneidad de preferencias}: ¿Cómo varían los determinantes de satisfacción según el perfil del usuario? (familias vs. profesionales jóvenes vs. adultos mayores vs. inversionistas)
\end{enumerate}

\subsubsection{Preguntas Secundarias}

\begin{itemize}
    \item \textbf{Autocorrelación espacial}: ¿Existe dependencia espacial en la satisfacción residencial? Es decir, ¿viviendas cercanas tienden a tener niveles similares de satisfacción debido a características compartidas del entorno?
    
    \item \textbf{Importancia relativa de factores}: ¿Cuáles son las 10 variables más determinantes de satisfacción? ¿Son características internas (superficie, dormitorios) o factores espaciales externos (distancia a metro, áreas verdes)?
    
    \item \textbf{Gradientes urbanos}: ¿Existen patrones geográficos sistemáticos? ¿La satisfacción aumenta/disminuye con la distancia al centro urbano o sigue otros patrones espaciales?
    
    \item \textbf{Equidad territorial}: ¿Qué comunas ofrecen mejor relación calidad-precio? ¿Existen zonas con alta satisfacción predicha pero precios relativamente accesibles?
    
    \item \textbf{Capacidad predictiva}: ¿Es posible predecir con alta precisión (R² > 0.85) la satisfacción residencial utilizando únicamente características observables y datos geográficos públicos?
    
    \item \textbf{Accesibilidad diferenciada}: ¿Qué dimensiones de accesibilidad (transporte, educación, salud, seguridad, áreas verdes) tienen mayor peso en la satisfacción? ¿Este peso varía según el tipo de propiedad (casa vs. departamento)?
\end{itemize}

\subsubsection{Hipótesis de Trabajo}

Basándonos en la literatura de geografía urbana y economía espacial, planteamos:


latex\begin{itemize}
    \item \textbf{H1}: La satisfacción residencial está significativamente influenciada por factores espaciales externos, más allá del precio y características internas de la vivienda.
    
    \item \textbf{H2}: Existe autocorrelación espacial positiva en la satisfacción residencial (Moran's $I > 0$), indicando clustering de zonas de alta/baja satisfacción.
    
    \item \textbf{H3}: La importancia relativa de factores varía sustancialmente entre perfiles de usuario, siendo transporte crítico para profesionales jóvenes, educación para familias, y salud para adultos mayores.
    
    \item \textbf{H4}: Modelos de Machine Learning que integran características espaciales superan en capacidad predictiva ($R^2 > +10\%$) a modelos que solo consideran atributos de la propiedad.
\end{itemize}


\subsection{Alcance y Limitaciones}

\subsubsection{Alcance Geográfico}
Este proyecto considera información geoespacial para las siguientes cuatro comunas de la Región Metropolitana de Santiago:

\begin{itemize}
    \item Ñuñoa
    \item La Reina
    \item Santiago Centro
    \item Estación Central
\end{itemize}

\subsubsection{Datos de Propiedades}
El análisis contempla un total de \textbf{7,702 propiedades} en venta, distribuidas en:
\begin{itemize}
    \item 5,135 departamentos
    \item 2,567 casas
\end{itemize}

Con las siguientes características:
\begin{itemize}
    \item Rango de precios: 800 - 73,000 UF
    \item Superficie: 16 - 900 m²
    \item Fuente: Portal Inmobiliario (web scraping) + Google Maps API (geocodificación)
\end{itemize}

\subsubsection{Variables Analizadas}

\textbf{Variables de la propiedad:}
\begin{itemize}
    \item Precio total (UF) y precio por metro cuadrado (UF/m²)
    \item Superficie útil (m²)
    \item Número de dormitorios y baños
    \item Tipo de propiedad (casa o departamento)
    \item Coordenadas geográficas (latitud/longitud)
\end{itemize}

\textbf{Variables espaciales (distancias y densidades):}
\begin{itemize}
    \item \textbf{Transporte:} Distancia a estaciones de metro y estaciones de carga eléctrica
    \item \textbf{Educación:} Distancia a colegios, jardines infantiles y universidades
    \item \textbf{Salud:} Distancia a hospitales, consultorios, farmacias y clínicas
    \item \textbf{Seguridad:} Distancia a unidades de la PDI, cuarteles y bomberos
    \item \textbf{Comercio:} Distancia a tiendas y servicios comerciales
    \item \textbf{Áreas verdes y recreación:} Distancia a parques, plazas y espacios de ocio
    \item \textbf{Densidades:} Cantidad de servicios por km² en radios de 300m, 600m y 1000m
\end{itemize}

\subsubsection{Limitaciones del Proyecto}

El modelo de satisfacción residencial presenta las siguientes limitaciones que deben considerarse al interpretar los resultados:

\begin{enumerate}
    \item \textbf{Cobertura geográfica limitada:} Solo se analizan 4 comunas de la Región Metropolitana, por lo que los resultados no son generalizables a otras comunas o regiones de Chile.
    
    \item \textbf{Antigüedad y estado de conservación:} No se dispone de información sobre el año de construcción ni el estado actual de las propiedades.
    
    \item \textbf{Características del edificio:} No se consideran factores como presencia de ascensor, piso del departamento, vista, orientación, estacionamiento o bodega.
    
    \item \textbf{Datos estáticos:} Los datos corresponden a un período específico (2025) y no capturan la evolución temporal del mercado inmobiliario.
    
    \item \textbf{Factores socioeconómicos:} No se incluyen variables detalladas del nivel socioeconómico del vecindario más allá de la comuna.
    
    \item \textbf{Calidad de construcción:} No existe información disponible sobre la calidad de los materiales o terminaciones.
    
    \item \textbf{Solo propiedades en venta:} El análisis no incluye propiedades en arriendo, limitando el alcance a compradores potenciales.
    
    \item \textbf{Gastos adicionales:} No se dispone de información sobre gastos comunes, contribuciones o costos de mantenimiento.
\end{enumerate}
